\section{Introducción}

La Interacción Humano-Computadora (IHC) estudia la relación entre las personas y los sistemas computacionales, con el propósito de comprender cómo los usuarios interactúan con la tecnología y cómo deben diseñarse las interfaces para que esa interacción resulte efectiva, eficiente, comprensible y satisfactoria. Desde esta perspectiva, un sistema no puede considerarse adecuado únicamente por su correcto funcionamiento técnico: si el usuario no logra entenderlo, usarlo o completar sus tareas sin un esfuerzo excesivo, el sistema falla en la dimensión que realmente importa, la de la interacción \autocite{nielsen1994,norman2002}.

En este contexto, no solo es relevante el aspecto visual de una interfaz, sino también la forma en que la información se organiza, estructura y presenta. La Arquitectura de la Información adquiere aquí un papel clave, ya que permite que los contenidos sean claros, accesibles y comprensibles, facilitando la navegación y reduciendo la carga cognitiva que el usuario debe asumir durante el uso del sistema \autocite{rosenfeld2015}.

El presente proyecto aplica estos principios al análisis y rediseño del Sistema de Información San Simón (WebSIS) de la Universidad Mayor de San Simón, plataforma que constituye el único canal institucional habilitado para que los estudiantes realicen trámites académicos como la inscripción a materias, la consulta de notas y horarios y la revisión de su situación académica. Al ser de uso obligatorio, sus deficiencias de interacción afectan de manera transversal a toda la población estudiantil, con mayor impacto durante los periodos críticos de inscripción y publicación de calificaciones.

A lo largo del trabajo se integran metodologías de investigación de usuarios, principios de usabilidad, factores cognitivos, arquitectura de la información, metodologías centradas en el usuario y técnicas de evaluación, con el fin de proponer una solución que mejore la experiencia de uso. El documento recorre el proceso completo de Diseño Centrado en el Usuario: parte de la definición del problema y la investigación de usuarios reales, continúa con la arquitectura de la información y el diseño de la interfaz, materializa la propuesta en un prototipo funcional y culmina con su evaluación de usabilidad. De esta manera, el proyecto busca no solo fortalecer habilidades técnicas, sino también fundamentar cada decisión de diseño en evidencia obtenida directamente de los usuarios.
